\begin{textsample}{POS Dim 4 – ms – Score 10.00 – ms/ms\_em\_3.txt}  \label{ex:f4_pos_154}
2.
A \textbf{TERRITORIALIDADE} DE MATO GROSSO DO \textbf{SUL} O \textbf{Estado} de Mato Grosso do Sul tem uma \textbf{população} estimada de 2.748.023 habitantes, segundo dados do Instituto Brasileiro de Geografia e Estatística - IBGE/2018, com uma área territorial de 357.145,531 km² e 79 municípios, sendo o 6º \textbf{estado} brasileiro no que diz respeito à extensão territorial.
Está localizado na região Centro-Oeste e seus limites geográficos são : Goiás ( \textbf{nordeste} ), Minas Gerais ( leste ), Mato Grosso ( \textbf{norte} ), Paraná ( \textbf{sul} ), São Paulo ( sudeste ), Bolívia ( oeste ), Paraguai ( \textbf{sul} e oeste ).
A maior parte do \textbf{Estado} é formada pela planície do Pantanal e pelo bioma do Cerrado, com clima tropical e solo fértil formado por terra roxa, o que acabou definindo o seu perfil socioeconômico voltado para a agricultura e pecuária, principalmente.
O \textbf{Estado} foi criado em 1977, em um momento singular da história brasileira, período do regime militar e também de desenvolvimento do Oeste, iniciado no governo de Getúlio Vargas que lançou um novo projeto para intensificar a ocupação do Centro-Oeste, conhecido como a “ Marcha para o Oeste ”, na década de 1940.
Segundo o Censo do IBGE do ano 2010, o \textbf{Estado} de Mato Grosso do Sul possui a segunda maior \textbf{população} \textbf{indígena} do país, com cerca de 77.025 \textbf{indígenas}, distribuídos em 75 aldeias ( dados DSEI/Funasa-MS - 2015 ), localizadas em 27 municípios, representados por oito etnias oficiais : Atikum, Guató, Guarani, Kaiowá, Kinikinau, Kadwéu, Ofaié e Terena.
Cada um desses grupos étnicos é um conjunto cultural único, com suas tradições, manifestações culturais e línguas.
São quatro troncos linguísticos encontrados no \textbf{Estado} de Mato Grosso do Sul : Tupi-Guarani, Macro Jê, Arúak e Guaikuru.
As \textbf{migrações} de contingentes oriundos dos Estados de Minas Gerais, Rio Grande do Sul, Paraná e São Paulo, as imigrações de países, como Alemanha, Espanha, Itália, Japão, Paraguai, Bolívia, Portugal, Síria e Líbano foram fundamentais para o povoamento de Mato Grosso do Sul e marcam a fisionomia dessa região.
O \textbf{Estado} também recebeu ciclos \textbf{migratórios} de quilombolas, remanescentes de Minas Gerais e Goiás, os quais também foram responsáveis pela formação socioeconômica ; atualmente, os quilombolas estão distribuídos em vinte e duas comunidades e mantêm, ainda, suas práticas culturais.
Segundo informações do Programa das Nações Unidas para o Desenvolvimento ( PNUD ) e do Instituto de Pesquisa Econômica Aplicada ( IPEA ) de 2010, o \textbf{Estado} de Mato Grosso do Sul apresenta um Índice de Desenvolvimento Humano ( IDH ) de 0,729, ocupando o décimo lugar no ranking dos \textbf{estados} brasileiros, o que situa essa Unidade Federativa ( UF ) na faixa de Desenvolvimento Humano Alto.
O IDH foi calculado com base nos dados do Censo Demográfico de 2010, do IBGE, e considera questões como longevidade, renda e educação.
Com relação à educação, as redes públicas e instituições privadas de ensino do \textbf{Estado} de Mato Grosso do Sul, na etapa do Ensino Fundamental, registraram 397.032 matrículas no ano 2019 ( INEP, 2020 ) ; desse quantitativo, 98,1 \% abrange a \textbf{população} na faixa etária de 6 a 14 anos, conforme o Anuário Brasileiro da Educação Básica 2020.
Na etapa do Ensino Médio, 103.482 estudantes foram matriculados no ano 2019 ( INEP, 2020 ) ; desse quantitativo, 67,4 \% abrange a \textbf{população} na faixa etária de 15 a 17 anos, conforme o supramencionado Anuário.
Fonte : Própria ( 2020 ).

% matched lemmas: estado, indígena, migratório, migração, nordeste, norte, população, sul, territorialidade
\end{textsample}
